\scanpage{128}
% line head | text | 77
\mspage{77}
% line 1 | text | carré cocartésien
carré cocartésien
% line 1–2 | diagram | intersection, U1, U1 prime, X1
\diagram{\begin{tikzcd}[ampersand replacement=\&,row sep=small,column sep=small]
U_1\cap U'_1\arrow[r]\arrow[d] \& U_1\arrow[d]\\
U'_1\arrow[r] \& X_1
\end{tikzcd}}
% line 1–4 | text | du type envisagé
du type envisagé dans $W(7)$ ($X_1$ réunion de deux ouverts $U_1,U'_1$) dans un
autre (\inserted{d'}indices $2$). Considérons les conditions suivantes

% line 5 | text | a) f on the intersection
a) $f_{U_1\cap U'_1}\in W$

% line 6 | text | b) f on the opens
b) $f_{U_1}$ et $f_{U'_1}$ sont $\in W$

% line 7 | text | c) f
c) $f\in W$.

% line 8–11 | text | Alors deux parmi ces conditions
Alors deux parmi ces conditions impliquent la troisième.
En d'autres termes, quand l'une des \unsure{ces|les} conditions est satisfaite, les deux autres sont équivalentes
entre elles.

% line 12–14 | text | NB Cet axiome
\uline{NB} Cet axiome est équivalent à l'axiome dual,
\deleted{\unsure{dans}} en termes d'objets $X_1,X_2$ représentés \inserted{chacun} comme
réunion de deux fermés.

% line 15–18 | text | Corollaire
\uline{Corollaire} \deleted{\unsure{Comme|Comment}} \inserted{Supposons $W(8)$ satisfait. Alors} la conclusion \unsure{de|à} $W(7\mathrm{bis})$ (axiome
qu'on n'a pas à utiliser ici), sous sa forme complétée
(cf le NB qui le suit) est valable: deux parmi les
conditions a), b), c) impliquent la troisième.

% line 19–21 | text | Il suffit de regarder
Il suffit de regarder l'unique morphisme
du diagramme $(\unsure{\alpha|a})$ p.~76, dans le diagramme
trivial
% line 21 | diagram | four terminal categories e
\diagram{\begin{tikzcd}[ampersand replacement=\&,row sep=tiny,column sep=tiny]
\& e\arrow[dl]\arrow[dr] \&\\
e\arrow[dr] \& \& e\arrow[dl]\\
\& e \&
\end{tikzcd}}
% line 21–24 | text | correspondant au cas où
, correspondant au cas où
$X$ est la catégorie finale $e$, \inserted{et} $U',U''$ égaux à $X$. Alors le
corollaire n'est autre que $W(8)$ appliqué à ce
cas particulier.

% line 25–27 | text | Cet argument fait sentir
Cet argument fait sentir que $W(\unsure{7|8})$ est \inserted{sans doute} plus
fort que $W(7\mathrm{bis})$. Incité par ce sentiment, je \unsure{me|m'}
aperçois avec plaisir qu'il est fondé:
% archive sheet number: 128
